\begin{abstract}
Инженерия требований критически зависит от своевременного выявления логических противоречий. В исследовании рассматривается проблема оценки способности больших языковых моделей к многошаговому логическому выводу в задачах анализа текстов технических заданий. Существующие эталонные наборы данных обладают рядом недостатков, включая контаминацию обучающих выборок и отсутствие строгого контроля глубины логических связей. В данной статье предлагается методология синтетического тестирования, генерирующая текстовые наборы данных с контролируемым уровнем информационного шума и заданным порядком конфликта. Для математической верификации минимальных невыполнимых подмножеств и предотвращения поверхностных эвристик применяется SMT-решатель Z3. Введен метод структурного доказательства уникальности конфликта, позволяющий верифицировать наборы данных за линейное время. Архитектура программного комплекса включает изолированную среду для нейросимвольного исполнения кода. Предварительные эксперименты на модели класса 14B параметров подтверждают существенное превосходство нейросимвольного контура над прямым текстовым логическим выводом, вскрывая уязвимость последнего к информационному шуму.
\end{abstract}

\section{Введение: Проблема инженерии требований}
Спецификации требований служат фундаментальной основой для проектирования любых программных и инженерных систем. Значение корректной постановки требований невозможно переоценить, поскольку технически безупречная реализация ошибочно сформулированных задач приводит к созданию неработоспособного продукта. В системной инженерии традиционно выделяют несколько типов требований: функциональные, требования к производительности, вычислительным ресурсам и взаимодействию компонентов.

По мере роста сложности систем неизбежно возникает риск появления конфликтующих требований. Конфликт требований — это ситуация, при которой их одновременное выполнение приводит к существенному ограничению пространства допустимых решений, вплоть до полного отсутствия такового. Специфика проблемы заключается в том, что требования могут быть выполнимыми по отдельности, но становиться абсолютно несовместимыми при совместном рассмотрении \cite{lamsweerde1998}.

Ошибки, заложенные на этапе составления технического задания, приводят к экспоненциальному увеличению финансовых и трудовых затрат, если скрытые конфликты проявляются лишь на поздних этапах жизненного цикла. Поэтому практическая значимость автоматического выявления конфликтов крайне высока для разработки аппаратно-программных комплексов, банковских и киберфизических систем. В связи с этим постановка задачи выявления конфликтующих требований выделяется в самостоятельную прикладную задачу, которая сегодня может решаться не только традиционными формальными методами, но и с применением больших языковых моделей.

\section{Предшествующие исследования}

\subsection{Конфликты требований и порядок конфликта}
Традиционные методы выявления конфликтов исторически базируются на попарном сравнении требований. Однако интуитивная идея попарного анализа обладает критическим ограничением: конфликт может отсутствовать во всех парах, но возникать при совместном рассмотрении трех и более условий. 

В работах А. Саладо и Р. Нилчиани \cite{salado2016} математически доказана несостоятельность исключительно попарного подхода и введена концепция порядка конфликта. Данная метрика связана с размером подмножества требований, которые необходимо рассмотреть одновременно для обнаружения противоречия. Анализ первого порядка охватывает подмножества из двух требований, в то время как более высокий порядок отражает конфликты, возникающие только при взаимодействии большего числа спецификаций. Таким образом, порядок конфликта определяет фундаментальный компромисс между полнотой анализа и его вычислительной стоимостью, требуя рассматривать набор требований как единое целое.

\subsection{Языковые модели для логического вывода}
Параллельно наблюдается стремительный рост применения больших языковых моделей для анализа технических документов и логического вывода. Существуют исследования по непосредственному обнаружению противоречий в естественно-языковых требованиях с использованием передовых коммерческих моделей. 

Однако исследования пределов композиционности архитектуры трансформеров \cite{dziri2023} показывают, что языковые модели имеют существенные ограничения при выполнении многошагового вывода: с ростом графа логической задачи качество работы моделей демонстрирует экспоненциальное падение. Прямой текстовый подход также сталкивается с серьезными трудностями при необходимости не просто ответить на вопрос о наличии конфликта, а восстановить точное противоречивое ядро.

\subsection{Нейросимвольные подходы и формальные решатели}
Стремление компенсировать указанные недостатки привело к появлению нейросимвольных подходов. В современных работах \cite{satlm2023, linc2023} языковая модель используется для генерации формального представления, а окончательное решение принимает внешний детерминированный решатель. Это позволяет разделить задачу понимания естественного языка и задачу строгой логической проверки.

В качестве инструмента автоматического вывода традиционно рассматриваются SMT-решатели. При выборе инструмента для данного исследования рассматривались различные альтернативы, однако выбор был сделан в пользу Z3 \cite{demoura2008} благодаря возможности реализации единого программного контура. Использование встроенных методов отслеживания утверждений в Z3 позволяет непосредственно связывать формализованные требования с их текстовыми идентификаторами, обеспечивая извлечение минимального невыполнимого подмножества без дополнительных эвристик.

\section{Исследовательская ниша}
Несмотря на существование наборов данных для оценки логики языковых моделей, в научной литературе недостаточно изучено влияние структурной сложности противоречия на качество автоматического вывода. В частности, открытыми остаются вопросы: как изменяется качество обнаружения противоречий при увеличении порядка конфликта, и каково влияние количества дополнительных непротиворечивых требований, создающих контекстный шум. Также требует систематической оценки сравнительная эффективность прямого текстового подхода и нейросимвольного подхода с формальной математической проверкой. Это формирует основной исследовательский вопрос настоящей работы.

\section{Терминология и формальная рамка}
Основным научным термином данной статьи является «порядок конфликта». В рамках разрабатываемого инструментария формальным представлением порядка конфликта выступает размер минимального невыполнимого подмножества. 

Минимальное невыполнимое подмножество  — это такое множество требований, которое удовлетворяет двум условиям:
\begin{enumerate}
    \item Выполнение всего множества требований одновременно невозможно.
    \item Удаление любого единичного требования из этого множества устраняет противоречие (условие минимальности).
\end{enumerate}

Размер данного подмножества выступает количественной характеристикой сложности рассматриваемого конфликта. Следует отметить установление соответствия с терминологией предшественников: в ранних работах порядок формально задавался через подмножество размера $n+1$. В настоящей статье под термином «порядок конфликта» понимается непосредственно количество требований, образующих противоречивое ядро.

\section{Разработанный тестовый набор и алгоритм генерации}
Целью исследования является разработка открытого инструментария для систематической оценки способности языковых моделей обнаруживать противоречия. Разработанный программный комплекс позволяет генерировать тестовые случаи с возможностью независимого управления порядком конфликта и общим количеством требований за счет внедрения непротиворечивого контекстного шума. 

Архитектура генератора основана на парном формировании тестовых случаев. Для каждого сценария создается пара: выполнимый и невыполнимый варианты. Ключевой особенностью является использование одинакового начального значения генератора псевдослучайных чисел для пары. Это обеспечивает строгую синхронизацию отображения абстрактных переменных на предикаты естественного языка, сохраняет одинаковую структуру шумовых требований и равный суммарный объем текста. Идентификаторы требований переназначаются на нейтральные, скрывая от анализирующей модели разницу между смысловым ядром и информационным шумом.

\section{Виды логических паттернов и стратегии конфликтов}
Топологическая классификация стратегий генерации охватывает семь базовых классов связей. Паттерны обобщены для произвольного количества переменных $n$. Пусть $V = \{x_1, x_2, \dots, x_m\}$ — множество пропозициональных переменных. Для наглядности механизмы трансляции формальной логики в естественный язык проиллюстрированы на примерах.

\begin{enumerate}
    \item \textbf{Прямое противоречие.} Базовый конфликт минимального порядка.
    \begin{equation} x_1 \land \neg x_1 \models \bot \end{equation}
    
    \item \textbf{Транзитивная цепочка.} Последовательность логических следствий длины $n$. Порядок конфликта равен $n+1$.
    \begin{equation} x_1 \land \left( \bigwedge_{i=1}^{n-1} (x_i \rightarrow x_{i+1}) \right) \land \neg x_n \models \bot \end{equation}
    \textit{Пример для $n=3$:}
    \begin{itemize}
        \item [$R_1$:] Пользователь успешно прошел авторизацию в корпоративном приложении.
        \item [$R_2$:] Успешная авторизация в приложении инициирует открытие защищенной сессии.
        \item [$R_3$:] Открытие защищенной сессии автоматически требует активации протокола шифрования.
        \item [$R_4$:] Использование данного протокола шифрования строго запрещено настройками сетевого шлюза.
    \end{itemize}

    \item \textbf{Ветвление.} Одно начальное условие $x_0$ инициирует $n$ альтернативных следствий, каждое из которых противоречит терминальному состоянию $y$. Порядок конфликта равен $n+3$.
    \begin{equation} \left( x_0 \rightarrow \bigvee_{i=1}^{n} x_i \right) \land \left( \bigwedge_{i=1}^{n} (x_i \rightarrow \neg y) \right) \land x_0 \land y \models \bot \end{equation}
    
    \item \textbf{Конвергентный конфликт (Слияние).} Схождение $n$ путей к взаимно исключающим требованиям.
    \begin{equation} \left( \bigwedge_{i=1}^{n-1} (x_i \rightarrow y) \right) \land (x_n \rightarrow \neg y) \land \left( \bigwedge_{i=1}^{n} x_i \right) \models \bot \end{equation}

    \item \textbf{Покрытие вариантов.} Ограничение на мощность множества доступных ресурсов.
    \begin{equation} \left( \bigvee_{i=1}^{n} x_i \right) \land \left( \bigwedge_{i=1}^{n} \neg x_i \right) \models \bot \end{equation}

    \item \textbf{Идемпотентный конфликт.} Две противоречащие импликации, исходящие из одного антецедента (порядок $2$).
    
    \item \textbf{Математические и числовые отношения.} Построение невыполнимых ядер через транзитивность числовых операторов строгого неравенства и эквивалентности.
\end{enumerate}

\section{Математическая верификация и структурное доказательство уникальности}
Каждый сгенерированный тестовый случай проходит обязательную независимую проверку средствами SMT-решателя. Традиционная проверка минимальности и уникальности ядра путем полного перебора всех подмножеств обладает экспоненциальной вычислительной сложностью, что делает генерацию глубоких конфликтов (порядка $9$ и выше) неприемлемо долгой. 

В данной работе разработана и применена теорема структурной уникальности. Доказательство строится от противного на основе анализа графа зависимостей переменных. Пусть $M$ — множество целевых правил ядра, а $C$ — транзитивное замыкание по переменным множества $M$. Пусть $N$ — множество шумовых правил. Алгоритм генератора гарантирует, что целевые и шумовые правила оперируют строго непересекающимися множествами переменных ($V(C) \cap V(N) = \emptyset$). 

Следовательно, любое подмножество шумовых правил $N$ всегда выполнимо (являясь набором независимых булевых присваиваний). Любое невыполнимое подмножество объединения $C \cup N$ не может использовать правила из $N$ для создания конфликта, поскольку переменные изолированы. Таким образом, конфликт инкапсулирован исключительно внутри подмножества $C$. Поскольку каждая из предложенных стратегий генерации детерминировано строит ровно одну замкнутую цепочку внутри $C$, глобальное ядро является строго уникальным. Этот метод снижает сложность верификации уникальности до квадратичного времени от числа правил.

\section{Контроль лексических утечек}
Для исключения возможности определения класса тестового случая на основе поверхностных текстовых признаков формируются строго сбалансированные классы. Оценка утечки проводится посредством векторизации на основе частотности термов на уровне слов и символьных N-грамм с последующим обучением логистической регрессии в режиме стратифицированной групповой перекрестной проверки. 
Установлено жесткое математическое ограничение: среднее значение F1-меры классификатора не должно превышать 0,5. Для маскировки паттернов отрицания применяется единый стандартизированный лингвистический маркер («неверно, что»), что лишает языковую модель возможности ориентироваться на грамматические артефакты при классификации.

\section{Методология оценки и метрики}

\subsection{Прямой и нейросимвольный подходы}
Оценка проводится по двум независимым направлениям:
\begin{itemize}
    \item \textbf{Прямой вывод:} Модель получает текст и возвращает структурированный ответ, содержащий логическое рассуждение, бинарный статус и массив идентификаторов правил ядра.
    \item \textbf{Нейросимвольный вывод (Z3):} Модель выполняет роль транслятора, преобразуя текст требований в код на языке Python с использованием библиотеки Z3 и функции отслеживания утверждений. Код исполняется в изолированной песочнице с жестким ограничением времени выполнения и блокировкой системных модулей.
\end{itemize}

\subsection{Таксономия метрик и ошибок}
Первично оценивается правильность бинарной классификации (Accuracy). Оценка качества поиска ядра осуществляется вычислением точности, полноты и F1-меры восстановления подмножества, а также доли точных совпадений (Exact Match). 

Введен расширенный аппарат диагностики отказов, классифицирующий ошибки модели на следующие категории:
\begin{itemize}
    \item Ошибочная бинарная классификация.
    \item Неполное извлечение ядра (пропуск критических правил).
    \item Избыточное извлечение (включение шумовых правил в конфликт).
    \item Отказ от ответа (модель сообщает о невозможности логического вывода из-за недостатка данных).
    \item Инфраструктурные ошибки (синтаксические ошибки и исключения времени выполнения в сгенерированном коде).
\end{itemize}
Дополнительно регистрируются объем затраченных токенов и время задержки ответа.

\section{Анализ предварительных экспериментальных данных}
Для валидации методологии была проведена серия предварительных тестов на модели \texttt{Qwen2.5-Coder-14B-Instruct}. Эксперименты разделены на две группы: оценка устойчивости к информационному шуму при фиксированном ядре и оценка устойчивости к увеличению размера конфликта при фиксированном общем объеме текста.

\subsection{Влияние информационного шума}
При фиксированном порядке конфликта $K=3$ общее количество требований в тексте последовательно увеличивалось (10, 20 и 30 правил). 
Установлено, что эффективность прямого текстового подхода стремительно деградирует по мере роста числа отвлекающих факторов. Если при 10 правилах F1-мера восстановления ядра составляла 0,49, то при увеличении объема до 20 и 30 правил она снизилась до 0,29 и 0,35 соответственно. При этом метрика точного совпадения (Exact Match) упала до нуля: модель неизбежно пропускала важные правила или захватывала шумовые.

В противовес этому, нейросимвольный подход (Z3) продемонстрировал исключительную устойчивость. Значение F1-меры стабильно удерживалось на уровне 0,8 вне зависимости от количества шума, а метрика точного совпадения составила 0,8. Строгий математический аппарат SMT-решателя гарантированно отсекал избыточные спецификации.

\subsection{Влияние масштаба конфликта}
При фиксированном объеме документа (20 правил) порядок конфликта варьировался в значениях $K \in \{2, 5, 10\}$.
Прямой подход показал нестабильные результаты (F1-мера колебалась в диапазоне 0,52–0,82). Модель способна обнаружить противоречие в целом, однако метрика точного извлечения ядра оставалась на критически низком уровне (от 0,2 до 0,4). Многошаговые логические цепи на естественном языке приводят к потере фокуса трансформера.
При использовании Z3 метрика точного совпадения достигала 1,0 для конфликтов малого порядка ($K=2$) и оставалась на уровне 0,8 для глубоких конфликтов ($K=10$).

\subsection{Анализ вычислительных затрат и ошибок генерации}
Главным ограничением нейросимвольного подхода выступает нестабильность генерации программного кода и высокие вычислительные накладные расходы. В тестах с большим числом правил модель начинала допускать инфраструктурные ошибки (от 10\% до 30\% тестовых случаев завершались аварийно). Основной причиной сбоев стала попытка языковой модели применить недопустимые типы данных, например, использование оператора математического сравнения к объектам логического типа (BoolRef).

Кроме того, трансляция в код требует генерации в среднем в 1,5–2 раза большего количества токенов (до 2300 токенов на запрос против 1400 в прямом подходе), что приводит к пропорциональному увеличению времени задержки и финансовых затрат на вычисления.

\section{Заключение}
В статье представлена комплексная методология и инструментарий для оценки способностей языковых моделей выявлять скрытые структурные конфликты в системных требованиях. Применение структурного анализа графа зависимостей позволило решить проблему вычислительной сложности верификации тестовых наборов. 

Предварительные эксперименты уверенно подтверждают гипотезу о том, что для задач системной инженерии прямая генерация текста не обеспечивает необходимой надежности и точности локализации ошибок. Нейросимвольный подход выступает мощным средством преодоления информационной перегрузки, гарантируя извлечение математически строгих минимальных невыполнимых ядер, несмотря на повышенные затраты вычислительных ресурсов и требовательность к качеству кодогенерации применяемой языковой модели. Планируемые дальнейшие исследования будут направлены на масштабное сравнительное тестирование флагманских моделей различных архитектур и внедрение механизмов многоагентной автоматической коррекции сгенерированного кода.